\chapter*{Bài tập bổ sung}

\exercise Cho $(x; y; z) \in \mathbb{Z}^3$. Chứng minh rằng $5x^3 + 11y^3 + 13z^3 = 0 \implies x = y = z = 0$.

\exercise Chứng minh rằng \begin{equation*}
   \begin{cases}
      A_1 \implies A_2 \\
      A_2 \implies A_3 \\
      \vdots \\
      A_{n-1} \implies A_n \\
      A_n \implies A_1
   \end{cases} \iff \forall (i; j) \in \llbracket 1; n \rrbracket \times \llbracket 1; n \rrbracket, (A_i \iff A_j).
\end{equation*}

\exercise Chứng minh rằng có vô hạn số nguyên tố có dạng $4n - 1$ với $n \in \mathbb{N}$. Làm tương tự với $6n - 1$.

\exercise Giải thích tại sao  câu ``Mọi luật đều có ngoại lệ.'' nếu là mệnh đề thì là mệnh đề sai.

\exercise Cho $A$ và $B$ là hai tập hợp. Với mỗi phần, tìm tất cả các tập hợp $x$ sao cho khẳng định được cho là đúng:
\begin{multicols}{2}
   \begin{enumerate}
      \item $x \cup A = B $;
      \item $x \cap A = B $;
      \item $x \setminus A = B$;
      \item $A \setminus x = B$;
      \item $x \vartriangle A = B$.
   \end{enumerate}
\end{multicols}

\exercise Cho các tập hợp $x$, $y$, $z$. Xác định hai ánh xạ $g: x \to z$ và $h: y \to z$. Chứng minh rằng, để tồn tại $f:x \to y$ sao cho có được biểu đồ giao hoán
\begin{tikzpicture}[baseline=(current bounding box.center)]
   \node (y) at (1, 1) {$y$};
   \node (x) at (0, 0) {$x$};
   \node (z) at (2, 0) {$z$};
   
   \draw[guide arrow] (x) -- (y) node[midway, above left] {$f$};
   \draw[guide arrow] (y) -- (z) node[midway, above right] {$h$};
   \draw[guide arrow] (x) -- (z) node[midway, below] {$g$};
\end{tikzpicture}, điều kiện cần và đủ là \begin{equation*}
   \forall a \in x, \exists b \in y, g(a) = h(b).
\end{equation*}

% \solution

% Giả sử tồn tại $f: x\to y$ thỏa mãn biểu đồ giao hoán. Khi này, chúng ta có \begin{equation*}
%    \forall a \in x, g(a) = h(f(a)).
% \end{equation*}
% Do đó, với mỗi $a$ thuộc $x$, tồn tại $b$ thuộc $y$ ($b = f(a)$) sao cho $g(a) = h(b)$.

% Dưới một góc nhìn khác, nếu như chúng ta được cho $\forall a \in x, \exists b \in y, g(a) = h(b)$. Xây dựng quan hệ $F$ với định nghĩa sau:
% \begin{equation*}
%    \forall c \in x, \forall d \in y, (c\;F\;d \iff g(c) = h(d)).
% \end{equation*}

\exercise Cho hai tập hợp khác rỗng $B$ và $K$, và ánh xạ $a: B \to K$. Chứng minh rằng
\begin{enumerate}
   \item $a$ đơn ánh khi và chỉ khi tồn tại một toàn ánh $b$ từ $B$ đến $K$ sao cho $b \circ a = \mathds{1}_B$,
   \item $a$ toàn ánh khi và chỉ khi tồn tại một đơn ánh $c$ từ $B$ đến $K$ sao cho $a \circ c = \mathds{1}_A$.
\end{enumerate}

\exercise Cho tập hợp $B$ và $K$. Chứng minh rằng
\begin{center}
   tồn tại một đơn ánh từ $B$ vào $K \iff $ tồn tại một toàn ánh từ $K$ vào $B$. 
\end{center}

\exercise Cho $E$ là một tập hợp, $n \in \mathbb{Z}^+$, $A_0, \dots, A_n$ là những bộ phận của $E$ sao cho:
\[
    \varnothing = A_0 \subsetneq A_1 \subsetneq A_2 \dots \subsetneq A_n = E.
\]
Kí hiệu $B_1 = A_1 \setminus A_0, \dots, B_n = A_n \setminus A_{n-1}$. Chứng minh rằng $\{B_1; \dots; B_n\}$ là một phân hoạch của $E$.

\setcounter{exercise}{0}

